\begin{lemma}
\label{lemma-splitting-simplicial-groups}
Let $U$ be a simplicial abelian group.
Then $U$ has a splitting obtained by taking $N(U_0) = U_0$ and
for $m \geq 1$ taking
$$
N(U_m) = \bigcap\nolimits_{i = 0}^{m - 1} \Ker(d^m_i).
$$
Moreover, this splitting is functorial on the category
of simplicial abelian groups.
\end{lemma}

\begin{proof}
By induction on $n$ we will show that the choice of $N(U_m)$
in the lemma guarantees that (\ref{equation-splitting}) is
an isomorphism for $m \leq n$. This is clear for $n = 0$.
In the rest of this proof we are going to
drop the superscripts from the maps $d_i$ and $s_i$ in order
to improve readability. We will also repeatedly use the relations
from Remark \ref{remark-relations}.

\medskip\noindent
First we make a general remark.
For $0 \leq i \leq m$ and $z \in U_m$ we have
$d_i(s_i(z)) = z$. Hence we can write
any $x \in U_{m + 1}$ uniquely as
$x = x' + x''$ with $d_i(x') = 0$
and $x'' \in \Im(s_i)$
by taking $x' = (x - s_i(d_i(x)))$ and
$x'' = s_i(d_i(x))$. Moreover, the element
$z \in U_m$ such that $x'' = s_i(z)$
is unique because $s_i$ is injective.

\medskip\noindent
Here is a procedure for decomposing any $x \in U_{n + 1}$.
First, write $x = x_0 + s_0(z_0)$ with $d_0(x_0) = 0$.
Next, write $x_0 = x_1 + s_1(z_1)$ with
$d_1(x_1) = 0$. Continue like this to get
\begin{eqnarray*}
x & = & x_0 + s_0(z_0), \\
x_0 & = & x_1 + s_1(z_1), \\
x_1 & = & x_2 + s_2(z_2), \\
\ldots & \ldots & \ldots \\
x_{n - 1} & = & x_n + s_n(z_n)
\end{eqnarray*}
where $d_i(x_i) = 0$ for all $i = n, \ldots, 0$.
By our general remark above all of the $x_i$
and $z_i$ are determined uniquely by $x$.
We claim that
$x_i \in
\Ker(d_0) \cap
\Ker(d_1) \cap
\ldots \cap
\Ker(d_i)$
and
$z_i \in
\Ker(d_0) \cap
\ldots \cap
\Ker(d_{i - 1})$
for $i = n, \ldots, 0$.
Here and in the following
an empty intersection of kernels indicates
the whole space; i.e.,
the notation
$z_0 \in \Ker(d_0) \cap
\ldots \cap
\Ker(d_{i - 1})$
when $i = 0$ means $z_0 \in U_n$ with no restriction.

\medskip\noindent
We prove this by ascending induction on $i$.
It is clear for $i = 0$ by construction of $x_0$ and $z_0$.
Let us prove it for $0 < i \leq n$ assuming the result for $i - 1$.
First of all we have $d_i(x_i) = 0$ by construction.
So pick a $j$ with $0 \leq j < i$. We have
$d_j(x_{i - 1}) = 0$ by induction. Hence
$$
0 = d_j(x_{i - 1})
= d_j(x_i) + d_j(s_i(z_i))
= d_j(x_i) + s_{i - 1}(d_j(z_i)).
$$
The last equality by the relations of Remark \ref{remark-relations}.
These relations also imply that
$d_{i - 1}(d_j(x_i)) = d_j(d_i(x_i)) = 0$
because $d_i(x_i)= 0$ by construction.
Then the uniqueness in the general remark above shows the equality
$0 = x' + x'' = d_j(x_i) + s_{i - 1}(d_j(z_i))$
can only hold if both terms are zero. We conclude that
$d_j(x_i) = 0$ and by injectivity of $s_{i - 1}$ we also
conclude that $d_j(z_i) = 0$. This proves the claim.

\medskip\noindent
The claim implies we can uniquely write
$$
x = s_0(z_0) + s_1(z_1) + \ldots + s_n(z_n) + x_n
$$
with $x_n \in N(U_{n + 1})$ and
$z_i \in \Ker(d_0) \cap \ldots \cap \Ker(d_{i - 1})$.
We can reformulate this as saying that we have found a direct
sum decomposition
$$
U_{n + 1}
=
N(U_{n + 1})
\oplus
\bigoplus\nolimits_{i = 0}^{i = n}
s_i\Big(\Ker(d_0) \cap \ldots \cap \Ker(d_{i - 1})\Big)
$$
with the property that
$$
\Ker(d_0) \cap \ldots \cap \Ker(d_j)
=
N(U_{n + 1}) \oplus
\bigoplus\nolimits_{i = j + 1}^{i = n}
s_i\Big(\Ker(d_0) \cap \ldots \cap \Ker(d_{i - 1})\Big)
$$
for $j = 0, \ldots, n$.
The result follows from this statement as follows.
Each of the $z_i$ in the expression for $x$
can be written uniquely as
$$
z_i = s_i(z'_{i, i}) + \ldots + s_{n - 1}(z'_{i, n - 1}) + z_{i, 0}
$$
with $z_{i, 0} \in N(U_n)$ and
$z'_{i, j} \in \Ker(d_0) \cap \ldots \cap \Ker(d_{j - 1})$.
The first few steps in the decomposition of $z_i$ are zero because
$z_i$ already is in the kernel of $d_0, \ldots, d_i$.
This in turn uniquely gives
$$
x = x_n + s_0(z_{0, 0}) + s_1(z_{1, 0}) + \ldots + s_n(z_{n, 0}) +
\sum\nolimits_{0 \leq i \leq j \leq n - 1} s_i(s_j(z'_{i, j})).
$$
Continuing in this fashion we see that we in the end obtain
a decomposition of $x$ as a sum of terms
of the form
$$
s_{i_1} s_{i_2} \ldots s_{i_k} (z)
$$
with $0 \leq i_1 \leq i_2 \leq \ldots \leq i_k \leq n - k + 1$ and
$z \in N(U_{n + 1 - k})$. This is exactly the required
decomposition, because any surjective map $[n + 1] \to [n + 1 - k]$
can be uniquely expressed in the form
$$
\sigma^{n - k}_{i_k} \ldots \sigma^{n - 1}_{i_2} \sigma^n_{i_1}
$$
with $0 \leq i_1 \leq i_2 \leq \ldots \leq i_k \leq n - k + 1$.
\end{proof}
